The coverage planning application has four features that fall outside the
standard assumptions used above:
\begin{enumerate}
  \item \textbf{Degenerate diffusion.}
        The information quality states $q_i$ evolve deterministically;
        the diffusion matrix is block-diagonal with a zero block,
        $\Sigma = \operatorname{diag}(\Sigma_x,\sigma_b^2,\mathbf{0}_{m\times m})$.
        Deep BSDE methods require a generalized Feynman-Kac theorem for
        this case; PINN/DGM handles it directly.
  \item \textbf{Non-quadratic Hamiltonian.}
        The information gain rate $\alpha_i(x,p_s)$ is nonlinear in
        the sensing power $p_s$, so the $\sup_u[\cdot]$ in the HJB
        does not reduce to a closed-form quadratic minimizer.
        The first-order necessary conditions must be solved numerically
        at each collocation point, or an actor network must be added.
  \item \textbf{Ito correction cost.}
        The second-order term $\frac{1}{2}\operatorname{Tr}(\Sigma\nabla_\xi^2 V)$
        involves only the $d\times d$ position Hessian and one scalar
        (battery), regardless of $m$.  Second-order AD cost is
        therefore $O(d^2+1)$, not $O((d+1+m)^2)$.
  \item \textbf{Multiple solutions.}
        The infinite-horizon stochastic HJB~\eqref{eqn:ss-HJB} has multiple
        solutions.  The finite-horizon approximation of~\cite{fotiadis2025pinn}
        resolves this; horizon extension must be applied.
\end{enumerate}
Based on these considerations, PINN/DGM applied to the finite-horizon
stochastic HJB is the recommended primary method for the coverage problem.
